\textbf{结论名称：} 空集是任何集合的子集，空集是任何非空集合的真子集。

\textbf{适用条件：} 结论对任意集合 \(A\) 成立，其中 \(\emptyset\) 为空集，是集合论的基本对象。

\textbf{变量说明：} \(\emptyset\)：空集（不包含任何元素的集合）；\(A\)：所讨论的任意集合。

\textbf{核心公式：}
\[
\boxed{\emptyset \subseteq A \quad \text{且} \quad A \neq \emptyset \implies \emptyset \subsetneqq A}
\]

\textbf{使用场景：} 判断集合间的包含关系；讨论子集、真子集的存在性与计数；在集合证明题中作为默认成立的基本事实直接引用。

\textbf{简单例子：} 设 \(A = \{1, 2\}\)，则 \(A\) 非空。依结论：\(\emptyset \subseteq A\) 成立（空集不含元素，其所有元素自然都在 \(A\) 中）；同时 \(\emptyset \subsetneqq A\) 成立，因为存在 \(1 \in A\) 但 \(1 \notin \emptyset\)。若 \(A = \emptyset\)，则只有 \(\emptyset \subseteq \emptyset\)（集合是自身的子集），此时 \(\emptyset\) 不是真子集，这与“非空”的限制一致。